% Generated draft: requires mathematical and source review.
\documentclass[11pt]{article}
\usepackage[margin=25mm]{geometry}
\usepackage{fontspec}
\setmainfont{Latin Modern Roman}
\usepackage{amsmath,amssymb,mathtools,unicode-math}
\setmathfont{Latin Modern Math}
\usepackage{enumitem,needspace,xurl,hyperref}
\setlength{\parindent}{0pt}
\setlength{\parskip}{5pt}
\setlength{\emergencystretch}{3em}
\newcommand{\N}{\mathbb{N}}
\newcommand{\Z}{\mathbb{Z}}
\newcommand{\Q}{\mathbb{Q}}
\newcommand{\R}{\mathbb{R}}
\newcommand{\C}{\mathbb{C}}
\newcommand{\F}{\mathbb{F}}
\newcommand{\A}{\mathbb{A}}
\newcommand{\PP}{\mathbb{P}}
\newcommand{\E}{\mathbb{E}}
\newcommand{\eps}{\varepsilon}
\newcommand{\dd}{\,\mathrm{d}}
\begin{document}
\section*{UnsolvedMath ID: 30005715}
\textbf{Source problem code:} catalogue entry
\subsection{Definitions and mathematical statement}
The supplied problem concerns the sequence
\[
 t_N=\#\{\text{combinatorial isomorphism classes of triangulations of homology three-spheres with }N\text{ facets}\},\qquad N\geq 1.
\]
For a finite abstract simplicial complex of dimension three, a facet is a maximal simplex; in a pure three-dimensional triangulation every facet is a tetrahedron. Two complexes are combinatorially isomorphic if a bijection of their vertex sets carries simplices to simplices in both directions.

The source record does not specify homology coefficients or explicitly define the permitted triangulations. Under the possible convention that a triangulation is a finite simplicial complex whose realization is a closed connected topological three-manifold, an integral homology three-sphere would satisfy
\[
 H_i(|K|;\Z)\cong
 \begin{cases}
 \Z,&i=0,3,\\
 0,&i\ne 0,3.
 \end{cases}
\]
If homology over a field is intended, the corresponding condition instead uses that field as coefficients. These conventions are conditional explanations, not hypotheses established by the supplied record. It is also unresolved whether the intended class requires the manifold condition or includes complexes defined by homological conditions alone.

\textbf{Question.} With the intended class of homology three-sphere triangulations fixed, determine whether \(t_N\) grows exponentially or superexponentially as \(N\) increases.

A precise upper-bound interpretation asks whether
\[
 \exists C>1\;\exists N_0\in\N\;\forall N\geq N_0:
 \quad t_N\leq C^N,
\]
or instead
\[
 \forall C>1\;\forall N_0\in\N\;\exists N\geq N_0:
 \quad t_N>C^N.
\]
These are complementary alternatives, but the source does not specify that these are its intended meanings. In particular, the second alternative does not assert \(t_N^{1/N}\to\infty\), and the first does not assert a matching exponential lower bound. Any stronger intended growth requirement, and any restriction to attainable facet counts, require clarification.
\subsection{Short English statement}
Count triangulations of homology three-spheres with exactly \(N\) tetrahedral facets, identifying triangulations that are combinatorially isomorphic. Determine whether this count admits an exponential upper bound or has superexponential growth in the intended sense. The homology coefficients, permitted triangulations, and precise growth convention remain unresolved in the supplied record.
\subsection{Sources}
\begin{itemize}
\item Supplied local UnsolvedMath registry record, \texttt{data/unsolvedmath/problems.json}, permanent ID \(30005715\), title \emph{Growth of Triangulated Homology Three-Spheres}, exact source problem code \texttt{catalogue entry}.
\item Supplied report citation: \emph{Geometric, Algebraic and Topological Combinatorics} (2024), Oberwolfach Reports, pp.\ 3249--3318, DOI \texttt{10.4171/owr/2023/58}. Original Oberwolfach report: \url{https://doi.org/10.4171/owr/2023/58}. The registry gives proposed year 2023; the supplied citation labels the report 2024.
\item Attribution in the supplied original excerpt: Geva Yashfe, \emph{Counting triangulations of homology 3-spheres}, joint work with Karim Adiprasito, Marc Lackenby, Juan Souto, and, separately, Yuval Peled. The registry does not separately identify a proposer.
\item Supplied literature source: Adiprasito--Lackenby--Souto--Yashfe, 2025, \url{https://arxiv.org/abs/2509.06275}.
\item Literature assessment dated exactly \texttt{2026-08-22}: the registry labels the problem partially solved and states that the exponential-versus-superexponential question remains unanswered, while recent work constructs substantial families and transfers superexponential growth from higher-dimensional homology spheres to dimension three. This is supplied dated curation, not an independently verified assessment. The supplied background explicitly requests independent verification of time-sensitive status and cited claims.
\item The supplied registry marks the statement \texttt{corrected\_verified} and reports that it was reconstructed and checked against the official Oberwolfach report. This is a supplied verification claim, not verification performed for this draft.
\item Context in the original excerpt attributes exponential enumeration of triangulations of \(S^2\) with \(N\) triangles to Tutte, reference [9]. The supplied record does not contain the full bibliographic entry for [9].
\end{itemize}
\end{document}
